\subsection{高一力学方法专题}
\begin{tabularx}{\columnwidth}{XX}
  受力分析顺序 & 先重力，再弹力，再摩擦力，再其他外力 \\
  共点力平衡条件 & $\sum F_x=0,\ \sum F_y=0$ \\
  三力平衡图解特征 & 三力首尾相接形成闭合三角形 \\
  摩擦力方向判定 & 与相对运动（或趋势）方向相反 \\
  动态平衡常用工具 & 图解法 + 正交分解 + 余弦定理 \\
\end{tabularx}

\begin{formulanotebox}
\begin{itemize}[leftmargin=1.5em]
  \item 绳子只提供拉力且方向沿绳；杆可拉可压，固定杆可提供非沿杆约束力。
  \item 判断摩擦力方向先看“是否已相对滑动”，未滑动时再判断“相对运动趋势”。
  \item 多力平衡题优先判断是否可几何法（力三角形）快速求解，不行再做正交分解。
\end{itemize}
\end{formulanotebox}

\begin{keypointbox}
模型化思路比公式更重要：先识别连接关系（绳/杆/接触面），再做受力图，最后选“几何法或代数法”。
\end{keypointbox}

\begin{casetypebox}[title={\strut\textcolor{white}{\bfseries 常见题型：摩擦力方向与大小判定}}]
% TODO: 速度法、极端假设法、受力分析法综合判定题型待补充
\end{casetypebox}

\begin{casetypebox}[title={\strut\textcolor{white}{\bfseries 常见题型：绳杆模型}}]
% TODO: 活结/死结张力关系、自由杆/固定杆受力与平衡题型待补充
\end{casetypebox}

\begin{casetypebox}[title={\strut\textcolor{white}{\bfseries 常见题型：三个力动态平衡}}]
% TODO: 一个力方向变、一个力大小定时的动态范围与极值题型待补充
\end{casetypebox}

\begin{examplebox}
% TODO: 力学方法专题例题待补充
\end{examplebox}

\begin{examplesolutionbox}
% TODO: 对应例题解析待补充
\end{examplesolutionbox}
